\chapter{From \texttt{isaaclab train} to a checkpoint}
\label{ch:pipeline}

This chapter follows one training run from the command line to the saved
policy, naming the Isaac Lab function that performs each step. Nothing here is
specific to the tricycle except the names of our classes; every Isaac Lab
task on the develop branch goes through the same path.

\section{The command}

\begin{lstlisting}[style=shell]
.\isaaclab.bat train --rl_library rsl_rl --task Luna-Tricycle-Direct --num_envs 4
\end{lstlisting}

\begin{itemize}[leftmargin=1.6em]
  \item \code{isaaclab.bat} activates the venv's Python and runs the
        \code{isaaclab} console script, i.e.\ the function \code{cli()} in
        \srcpath{source/isaaclab/isaaclab/cli/__init__.py}.
  \item \code{train} is the sub-command. Others are \code{play},
        \code{zero_agent}, \code{random_agent}, \code{benchmark}.
  \item \code{--rl_library rsl_rl} picks the learning library backend.
  \item \code{--task Luna-Tricycle-Direct} is the Gymnasium id registered by
        \srcpath{tricycle/__init__.py}. There is no \code{-v0} suffix.
  \item \code{--num_envs 4} overrides \code{scene.num_envs} from the config.
  \item Optional: \code{--viz newton} opens the Newton OpenGL window,
        \code{--max_iterations N} overrides the runner config,
        \srcpath{physics=newton_mjwarp} is a Hydra ``typed selector'' that swaps
        the physics section for a preset (not needed here because our
        \code{__post_init__} builds its own \code{NewtonCfg}).
\end{itemize}

\section{The whole path at a glance}

\begin{figure}[H]
\centering
\begin{tikzpicture}[
  node distance=5mm and 8mm,
  box/.style={draw,rounded corners=2pt,align=left,font=\small,inner sep=5pt,minimum width=5.4cm,fill=white},
  ours/.style={box,draw=isaacblue,thick,fill=isaacblue!5},
  lab/.style={box,draw=black!60},
  arr/.style={-{Stealth[length=2mm]},thick,black!70},
  lbl/.style={font=\scriptsize\itshape,text=black!60,align=left}
]
\node[lab] (cli) {\textbf{isaaclab cli()}\\ \srcpath{isaaclab/cli/__init__.py}};
\node[lab,below=of cli] (ep) {\textbf{load entry points} \code{isaaclab.tasks}\\ imports \code{luna_hifi_tasks}};
\node[ours,below=of ep] (reg) {\textbf{gym.register} in \srcpath{tricycle/__init__.py}\\ id $\to$ env class, cfg class, agent cfg class};
\node[lab,below=of reg] (train) {\textbf{run\_train\_cli} $\to$ \srcpath{backends/train_rsl_rl.py}};
\node[lab,below=of train] (hydra) {\textbf{resolve configs} \srcpath{isaaclab_tasks/utils/hydra.py}\\ instantiate cfgs, \code{__post_init__}, presets, overrides};
\node[ours,below=of hydra] (cfgs) {\textbf{TricycleEnvCfg} + \textbf{TricyclePPORunnerCfg}\\ \code{--num_envs} applied here, after \code{__post_init__}};
\node[lab,below=of cfgs] (make) {\textbf{gym.make} $\to$ \srcpath{DirectRLEnv.__init__}\\ SimulationContext, InteractiveScene, sim.reset()};
\node[ours,below=of make] (env) {\textbf{TricycleEnv} \code{tricycle_env.py}\\ \code{_setup_scene}, joint ids, buffers};
\node[lab,below=of env] (wrap) {\textbf{RslRlVecEnvWrapper}\\ obs dict $\to$ TensorDict, \code{time_outs} in extras};
\node[lab,below=of wrap] (runner) {\textbf{OnPolicyRunner.learn()} (rsl\_rl)\\ rollout $\to$ GAE $\to$ PPO update, 300 iterations};
\node[lab,below=of runner] (ckpt) {\textbf{checkpoints} \srcpath{logs/rsl_rl/tricycle/<time>/model_*.pt}};
\foreach \a/\b in {cli/ep,ep/reg,reg/train,train/hydra,hydra/cfgs,cfgs/make,make/env,env/wrap,wrap/runner,runner/ckpt}
  \draw[arr] (\a) -- (\b);
\node[lbl,right=6mm of ep] {start-up time};
\node[lbl,right=6mm of hydra] {configuration time};
\node[lbl,right=6mm of make] {run time};
\end{tikzpicture}
\caption{The training path. Blue boxes are code in this repository; grey boxes are Isaac Lab or \code{rsl_rl}.}
\label{fig:path}
\end{figure}

\section{Step 1: the CLI finds our package}
\label{sec:step_cli}

\code{cli()} checks whether the first argument is one of its sub-commands. If
it is, it first calls \srcpath{_load_external_tasks()}:

\begin{lstlisting}[numbers=none]
_TASK_ENTRY_POINT_GROUP = "isaaclab.tasks"

def _load_external_tasks() -> None:
    for entry_point in importlib.metadata.entry_points(group=_TASK_ENTRY_POINT_GROUP):
        entry_point.load()
\end{lstlisting}

\begin{isaacbox}{entry points}
An \emph{entry point} is a line of metadata that a package writes into the
Python environment when it is installed. Isaac Lab defines the group name
\code{isaaclab.tasks}; any installed package that lists something under that
group gets imported by the CLI before the task name is looked up. Our
\code{pyproject.toml} lists \code{luna_hifi_tasks = "luna_hifi_tasks"}, so
\code{entry_point.load()} performs \code{import luna_hifi_tasks}.
Entry points are read from the installed metadata, not from the source file,
which is why editing \code{pyproject.toml} requires re-running
\code{pip install -e}.
\end{isaacbox}

Importing \code{luna_hifi_tasks} runs the chain described in
Section~\ref{sec:when}: the package \code{__init__} imports \code{tricycle},
whose \code{__init__} calls \code{gym.register}. After this, Gymnasium's global
registry maps the string \code{Luna-Tricycle-Direct} to three dotted paths:
the environment class, the environment config class and the \code{rsl_rl}
agent config class (Chapter~\ref{ch:task_init}).

\code{cli()} then calls \code{train()}, which imports
\srcpath{isaaclab_rl.entrypoints.run_train_cli}. That function reads
\code{--rl_library} and dispatches to
\srcpath{isaaclab_rl/entrypoints/backends/train_rsl_rl.py}.

\section{Step 2: configs are built and overridden}
\label{sec:step_cfg}

The backend uses the \code{hydra_task_config} decorator from
\srcpath{source/isaaclab_tasks/isaaclab_tasks/utils/hydra.py}. Its
\code{register_task} function does the following, in order:

\begin{enumerate}[leftmargin=1.8em]
  \item \code{load_cfg_from_registry(task, "env_cfg_entry_point")} looks up
        the registered string \srcpath{"...tricycle_env_cfg:TricycleEnvCfg"},
        imports that module and \emph{instantiates the class}. Instantiating a
        \code{configclass} runs its \code{__post_init__}, so at this moment
        our \srcpath{TricycleEnvCfg.__post_init__} builds the whole
        \code{NewtonCfg} (Chapter~\ref{ch:env_cfg}). The same is done for the
        agent config.
  \item Presets are collected and any \code{physics=NAME} or
        \code{presets=...} selectors on the command line are applied with
        \emph{replace} semantics: the whole config section is swapped, not
        merged.
  \item If this is a \code{play} run, \code{env_cfg.play_mode()} is called
        (Section~\ref{sec:play}).
  \item Remaining \code{env.xxx=value} / \code{agent.xxx=value} scalar
        overrides are set with \code{setattr}, and anything else is handed to
        Hydra proper.
\end{enumerate}

Separately, the backend's argument parser handles the classic flags. In
\srcpath{isaaclab_rl/entrypoints/common.py}, \code{apply_env_overrides} does

\begin{lstlisting}[numbers=none]
if getattr(args_cli, "num_envs", None) is not None:
    env_cfg.scene.num_envs = args_cli.num_envs
\end{lstlisting}

\begin{warnbox}{\code{--num_envs} arrives after \code{__post_init__}}
The config object is fully built, \code{__post_init__} included, before the
command-line \code{--num_envs} is written into it. Anything you computed
inside \code{__post_init__} from \code{self.scene.num_envs} used the
\emph{class default} (8), not the number you typed. This is why the MPM
sparse-grid capacities in \code{tricycle_env_cfg.py} are fixed numbers that
must be raised by hand when you train with more envs
(Chapter~\ref{ch:physics}).
\end{warnbox}

\section{Step 3: the environment is built}
\label{sec:step_env}

\code{gym.make("Luna-Tricycle-Direct", cfg=env_cfg, ...)} instantiates
\code{TricycleEnv}, whose \code{__init__} immediately calls
\srcpath{DirectRLEnv.__init__} (\srcpath{source/isaaclab/isaaclab/envs/direct_rl_env.py}).
That constructor does a great deal, in this order:

\begin{enumerate}[leftmargin=1.8em]
  \item \code{cfg.validate()}: every \code{MISSING} field must have been
        filled in (\code{decimation}, \code{episode_length_s}, \code{scene},
        \code{observation_space}, \code{action_space} are all required).
  \item Seeds torch/numpy if \code{cfg.seed} is set (the runner sets it).
  \item Creates the \code{SimulationContext} from \code{cfg.sim}. This is the
        object that owns the physics backend; because our
        \code{cfg.sim.physics} is a \code{NewtonCfg}, the Newton manager is
        chosen (Chapter~\ref{ch:physics}).
  \item \code{self.scene = InteractiveScene(self.cfg.scene)}: walks every
        field of \code{TricycleSceneCfg}, spawns the prims on the USD stage
        (ground plane, light, the car USD, the hidden slab, the soil particle
        spawner), and creates the Python asset objects (\code{Articulation},
        \code{RigidObject}, \code{MPMObject}) that will wrap them.
  \item \code{self._setup_scene()}: our hook. In this repository it only
        fetches the already-created assets by name.
  \item \code{self.sim.reset()}: starts physics. For Newton this is where the
        USD stage is read into a Newton \code{ModelBuilder}, replicated once
        per env, finalised into a \code{Model}, and where the coupled solver
        (rigid + MPM) is constructed. After this call, tensors like
        \code{car.data.root_pos_w} are live.
  \item \srcpath{self.scene.update(dt)} fills the data buffers once.
  \item Allocates the per-env bookkeeping tensors:
        \code{episode_length_buf}, \code{reset_terminated},
        \code{reset_time_outs}, \code{reset_buf}, and builds the Gymnasium
        spaces from \srcpath{observation_space} and \srcpath{action_space}
        (an \code{int} becomes a \code{Box(-inf, inf, (n,))}).
\end{enumerate}

Only when \code{super().__init__} returns does the rest of
\srcpath{TricycleEnv.__init__} run: looking up joint indices and allocating the
action and previous-position buffers (Chapter~\ref{ch:env}).

\begin{isaacbox}{InteractiveScene and asset objects}
\code{InteractiveScene} (\srcpath{source/isaaclab/isaaclab/scene/interactive_scene.py})
turns a scene \emph{config} into scene \emph{objects}. Each attribute of the
scene cfg whose type is an asset cfg (\code{ArticulationCfg},
\code{RigidObjectCfg}, \code{MPMObjectCfg}, \code{AssetBaseCfg}) becomes an
entry in the scene, retrievable with \code{scene["name"]}. The scene also owns
\code{env_origins}, an $(N,3)$ tensor with the world position of each env's
origin, spaced \code{env_spacing} metres apart on a grid.
\end{isaacbox}

\section{Step 4: the wrapper}

\code{rsl_rl} expects a specific interface (a \code{VecEnv}) rather than the
Gymnasium one. \code{RslRlVecEnvWrapper}
(\srcpath{source/isaaclab_rl/isaaclab_rl/rsl_rl/vecenv_wrapper.py}) adapts it:

\begin{itemize}[leftmargin=1.6em]
  \item \code{get_observations()} calls our \code{_get_observations()} and
        returns the dictionary as a \code{TensorDict}. Our env returns a
        dict with a single group, \code{"policy"}; \code{rsl_rl} maps that
        group to both the actor and the critic when no \code{obs_groups} are
        configured.
  \item \code{step(actions)} optionally clips actions to
        \code{clip_actions}, calls \code{DirectRLEnv.step}, merges
        \code{terminated | truncated} into a single \code{dones} tensor, and
        stores \code{truncated} as \code{extras["time_outs"]} so PPO can
        bootstrap (Section~\ref{sec:term_vs_timeout}).
  \item Exposes \code{num_envs}, \code{num_actions}, \code{episode_length_buf}
        and \code{device}.
\end{itemize}

\section{Step 5: the runner}

\code{train_rsl_rl.py} builds the log directory
\srcpath{logs/rsl_rl/<experiment_name>/<YYYY-MM-DD_HH-MM-SS>}, dumps both configs
into it as YAML (\srcpath{params/env.yaml}, \srcpath{params/agent.yaml}), then
creates the runner and starts learning:

\begin{lstlisting}[numbers=none]
runner = OnPolicyRunner(env, agent_cfg.to_dict(), log_dir=log_dir, device=agent_cfg.device)
runner.learn(num_learning_iterations=agent_cfg.max_iterations)
\end{lstlisting}

Inside \code{learn} (\srcpath{rsl_rl/runners/on_policy_runner.py}) every
iteration is:

\begin{lstlisting}[numbers=none]
for _ in range(num_steps_per_env):          # 50 in our config
    actions = self.alg.act(obs)               # sample from the Gaussian policy
    obs, rewards, dones, extras = self.env.step(actions)
    self.alg.process_env_step(obs, rewards, dones, extras)   # store + handle time_outs
self.alg.compute_returns(obs)                 # GAE
loss_dict = self.alg.update()                 # PPO epochs / mini-batches
# log, and every save_interval iterations: torch.save(..., f"model_{it}.pt")
\end{lstlisting}

\code{TensorBoard} event files are written next to the checkpoints, so
\code{tensorboard --logdir logs/rsl_rl/tricycle} shows the curves.

\section{One environment step, in detail}
\label{sec:step_detail}

\srcpath{DirectRLEnv.step(action)} is the heart of the whole system and the order
of operations inside it decides what your overridden methods may assume.
Figure~\ref{fig:step} shows it for our \code{decimation = 2}.

\begin{figure}[H]
\centering
\begin{tikzpicture}[
  node distance=3mm,
  st/.style={draw,rounded corners=2pt,font=\small\ttfamily,inner sep=4pt,minimum width=6.0cm,align=left,fill=white},
  ours/.style={st,draw=isaacblue,thick,fill=isaacblue!6},
  lab/.style={st,draw=black!60},
  note/.style={font=\scriptsize,text=black!65,align=left},
  arr/.style={-{Stealth[length=1.8mm]},black!70}
]
\node[ours] (pre) {\_pre\_physics\_step(action)};
\node[note,right=4mm of pre] {store clamped actions};
\node[lab,below=of pre] (loop) {for \_ in range(decimation):   \# 2 times};
\node[ours,below=of loop,xshift=6mm] (apply) {\_apply\_action()};
\node[note,right=4mm of apply] {write wheel / steer targets};
\node[lab,below=of apply] (wds) {scene.write\_data\_to\_sim()};
\node[lab,below=of wds] (simstep) {sim.step()   \# 10 ms of physics};
\node[note,right=4mm of simstep] {coupled rigid + MPM solve};
\node[lab,below=of simstep] (upd) {scene.update(dt)};
\node[note,right=4mm of upd] {refresh car.data.* tensors};
\node[lab,below=of upd,xshift=-6mm] (cnt) {episode\_length\_buf += 1};
\node[ours,below=of cnt] (dones) {terminated, time\_out = \_get\_dones()};
\node[ours,below=of dones] (rew) {reward = \_get\_rewards()};
\node[note,right=4mm of rew] {uses state \emph{before} any reset};
\node[lab,below=of rew] (ids) {reset\_env\_ids = nonzero(terminated | time\_out)};
\node[ours,below=of ids] (reset) {\_reset\_idx(reset\_env\_ids)};
\node[note,right=4mm of reset] {only for envs that ended};
\node[ours,below=of reset] (obs) {obs = \_get\_observations()};
\node[note,right=4mm of obs] {already sees the fresh state of reset envs};
\node[lab,below=of obs] (ret) {return obs, reward, terminated, time\_out, extras};
\draw[arr] (pre) -- (loop);
\draw[arr] (loop) -- (apply);
\draw[arr] (apply) -- (wds);
\draw[arr] (wds) -- (simstep);
\draw[arr] (simstep) -- (upd);
\draw[arr] (upd) -- (cnt);
\draw[arr] (cnt) -- (dones);
\draw[arr] (dones) -- (rew);
\draw[arr] (rew) -- (ids);
\draw[arr] (ids) -- (reset);
\draw[arr] (reset) -- (obs);
\draw[arr] (obs) -- (ret);
\begin{scope}[on background layer]
\node[draw=black!25,dashed,rounded corners,fit=(apply)(wds)(simstep)(upd),inner sep=3pt] {};
\end{scope}
\end{tikzpicture}
\caption{Order of operations in \code{DirectRLEnv.step}. Blue: methods implemented in \code{tricycle_env.py}. The dashed group runs \code{decimation} times per call.}
\label{fig:step}
\end{figure}

Three consequences that the tricycle code relies on:

\begin{itemize}[leftmargin=1.6em]
  \item \textbf{Rewards are computed before resets.} \code{_get_rewards} may
        compare the current car position with the position stored last step
        (\code{_prev_x}); at that moment no env has been teleported yet.
  \item \textbf{Observations are computed after resets.} The observation
        returned for an env that just ended already describes the
        \emph{new} episode. \code{rsl_rl} knows the episode boundary from
        \code{dones}, so this is the expected behaviour.
  \item \textbf{\code{_apply_action} runs every physics sub-step,}
        \code{_pre_physics_step} once per control step. Cheap work (writing
        targets) belongs in the former, decoding the action in the latter.
\end{itemize}

\begin{isaacbox}{the Newton fast path for decimation}
After \code{sim.reset()}, \code{DirectRLEnv} asks the physics manager whether
it can run the whole decimation loop itself
(\srcpath{physics_manager.set_decimation(...)} /
\srcpath{handles_decimation()}). When it can, the loop in
Figure~\ref{fig:step} collapses into a single \code{sim.step} that internally
performs \code{decimation} physics steps under one CUDA graph, and
\code{_apply_action()} is called once. The observable behaviour for our task
is identical because \code{_apply_action} writes the same targets each time.
\end{isaacbox}

\section{Play}
\label{sec:play}

\begin{lstlisting}[style=shell]
.\isaaclab.bat play --rl_library rsl_rl --task Luna-Tricycle-Direct --num_envs 1 --checkpoint latest --viz newton
\end{lstlisting}

\code{play} follows the same path with two differences. First,
\code{hydra_task_config(..., play_mode=True)} calls
\code{env_cfg.play_mode()} after the presets are resolved and before scalar
overrides, so our override in \code{TricycleEnvCfg} caps the env count and
switches on particle rendering (Chapter~\ref{ch:env_cfg}); \code{--num_envs 1}
then still wins because it is applied afterwards. Second, instead of learning,
\code{play_rsl_rl.py} resolves the checkpoint (\code{latest} means the newest
run directory under \srcpath{logs/rsl_rl/tricycle} and its highest-numbered
\code{model_*.pt}), loads it into the runner, takes
\srcpath{runner.get_inference_policy()} (the deterministic mean action), exports
the network to \srcpath{exported/policy.pt} and \code{policy.onnx}, and loops
\code{obs, _, dones, _ = env.step(policy(obs))} until the window closes.

\code{--viz newton} attaches the Newton OpenGL visualiser. It draws rigid
shapes by default; soil particles are drawn only because our
\code{play_mode()} sets \code{show_particles=True}.
